\documentclass[12pt]{article}
\usepackage[utf8]{inputenc}
\usepackage[T1]{fontenc}
\usepackage{hyperref}
\title{West Yugur Language Resources for K-12 Curriculum Development}
\author{Compiled by Claude Sonnet 5.5 (claude-sonnet-5-5) for Andr\'as Kornai}
\date{}
\begin{document}
\maketitle

\section{Introduction}
West Yugur (ISO 639-3 ybe; Glottocode west2402) is a Turkic language of China with about 4,600 speakers (2007 figure). It is classed as threatened. It is thought to be used as a first language mainly by older adults, and it is not known to be taught in schools. Little material exists, and nothing found is designed for schools.

\section{Primary Educational Resources}
No primer, textbook or school programme for West Yugur turned up in the searches. Anyone developing a curriculum would have to build the materials from scratch, working from the descriptive sources below.

\section{Linguistic and Reference Materials}
\begin{itemize}
\item Martina Erica Roos, \emph{The Western Yugur (Yellow Yugur) Language: Grammar, Texts, Vocabulary} (Leiden University PhD dissertation, 2000). It is cited in the Endangered Languages bibliography: \url{https://www.endangeredlanguages.com/node/37682/bibliography}
\item Endangered Languages Catalogue entry for the language: \url{https://endangeredlanguages.com/node/37682}
\item Wikipedia page on ISO 639 ybe: \url{https://wikipedia.com/wiki/ISO_639:ybe}
\item Wikipedia page on Yugur (a mirror), which mentions early word lists by Potanin (1893) and Mannerheim and recordings by Malov: \url{https://www.chise.org/ipns/en.wikipedia-on-ipfs.org/wiki/Yugur}
\end{itemize}

\section{Religious and Community Resources}
No religious or community-organisation resources turned up in the searches.

\section{Digital and Audio Resources}
\begin{itemize}
\item Living Dictionaries, which has a Western Yugur entry (only a single entry was seen in the results): \url{https://livingdictionaries.app/western-yugur/entry/qJeARECVqUpviesl0AXt}
\item Omniglot page on Western Yugur, which shows a Latin alphabet chart. It is unclear whether speakers use this alphabet: \url{https://www.omniglot.com/writing/westernyugur.htm}
\item Verbix page for West Yugur (a verb list): \url{https://docs.verbix.com/Languages/YugurWest}
\end{itemize}

\section{Curriculum Development Context}
The language has few speakers, and the sources suggest it is mostly spoken by older people. A native-language curriculum would depend on new documentation work with the community, starting from Roos's grammar and texts. No orthography in use in the community was confirmed.

\end{document}